% =============================================================================
% tex/example-tex/main.tex — the worksheet TEMPLATE, LaTeX flavour.
% (tex/example-mdx/main.mdx is the same skeleton in MDX. Pick one: a LaTeX
% sheet also becomes a PDF; an MDX sheet is a web page only.)
%
% Two jobs. (1) A skeleton in the shape of the Iliad Intensive module
% template (the Mega Doc's "Template" tab): copy this folder, paste your
% material into the matching sections, delete what you don't need.
% (2) A demo of every construct iliad.sty and the tex -> MDX pipeline
% support, each used where a real sheet would use it.
%
% Start a new sheet:      cp -r tex/example-tex tex/<your-slug>
% Then list <your-slug> under its day in schedule.yaml (the only other edit).
% Build + preview:        ./run.sh watch <your-slug>
% Reference:              docs/commands.md (every construct), docs/iliad-sty.md
%
% The `%` comments are instructions to you, the author: they never reach the
% web page (the converter strips them) but they DO travel in the downloadable
% .tex, so delete them as you fill each section in.
% =============================================================================

% ---- Metadata ---------------------------------------------------------------
% A YAML block in comments. LaTeX ignores it; the converter lifts it into the
% page frontmatter. Keys:
%   title:     overrides \title{} on the web (optional).
%   summary:   the sheet's OVERVIEW — one tight paragraph. The page header
%              shows it under the title and it doubles as the index blurb, so
%              it lives here, never as an "Overview" body section. Use the
%              YAML folded scalar (`>-`) with continuation lines indented two
%              spaces after the leading `% `.
%   slides:    a deck hosted elsewhere (Drive FOLDER url, for a PDF with no
%              source). A slides.tex / slides.typ in this folder is compiled
%              and listed automatically — see slides.tex.
%   unlisted:  maintainer flag — built, linked from nowhere. Remove it.
% NOT here: cluster / day. schedule.yaml owns where a sheet sits in the course
% and the build stamps both in; writing them here is a build error.
%--- iliad -------------------------------------------------------------------
% title: "Example Worksheet: The Kitchen Sink"
% summary: >-
%   A template worksheet in the shape of the module template, demonstrating
%   every construct the pipeline supports: this paragraph is the overview,
%   shown under the title and in the index.
% # slides: https://drive.google.com/drive/folders/<folderId>
% unlisted: true
%--- end ----------------------------------------------------------------------
\documentclass[11pt]{article}

% The ONE required package. It provides: exercise, solution, hint, proof,
% callout, remark, teachingnote, definition/theorem/lemma/proposition/
% corollary/fact/example, learningoutcomes, solutionsonly, pdfonly, \youtube,
% the \ifsolutions toggle, and loads hyperref+cleveref (configure hyperref
% with \hypersetup{...}; don't re-load it). PDF styling is plain by default;
% [boxes] gives the ready-made colour boxes matching the website.
% Local copy first (Overleaf), else the shared ../iliad.sty:
\IfFileExists{iliad.sty}{\usepackage[boxes]{iliad}}{\usepackage[boxes]{../iliad}}

% \solutionsfalse   % uncomment to hide all solutions from your own PDF

% ---- Free zone --------------------------------------------------------------
% Define whatever notation you like; load extra packages as needed.
% \newcommand macros become KaTeX macros automatically (avoid \mathchoice and
% optional-argument macros — the converter warns if it can't translate one).
\usepackage{tikz}
\usetikzlibrary{positioning}
\newcommand{\R}{\mathbb{R}}
\newcommand{\E}{\mathbb{E}}
\newcommand{\KL}[2]{D_{\mathrm{KL}}\!\left(#1 \,\middle\|\, #2\right)}
\DeclareMathOperator*{\argmax}{arg\,max}

% ---- Title and contributors -------------------------------------------------
% \author{} is the Contributors section: everyone who wrote or taught the
% material, with affiliation. The web byline reads "Name (Affiliation)".
\title{Example Worksheet: The Kitchen Sink}
\author{\authorname{Jane Doe}\\ \affiliation{Example University}
  \and \authorname{John Roe}\\ \affiliation{Example Institute}}
\date{\today}

\begin{document}
\maketitle

% Optional. LaTeX's ToC in the PDF; an in-page Contents list on the web.
\tableofcontents

% ---- Front matter -----------------------------------------------------------
% Every sheet opens in this fixed order (the build warns when one strays):
%   1. video embeds (optional)   2. Prerequisites   3. the learning-outcomes
%   box (= the module intent)  — then the content.
% The overview is NOT a body section: it is the summary: in the metadata.

% \youtube[Title]{VIDEO_ID} — the 11-character ID (the v= value of the watch
% URL), never the full URL. Embedded player on the web; a "Video:" line with
% the watch URL in the PDF. With no [Title] the build asks YouTube for it.
% A lecture recording goes here, first. A video illustrating one point may
% sit with that point instead.
\youtube{dQw4w9WgXcQ}

%===============================================================================
\section{Prerequisites}
\label{sec:prereqs}
% What should a participant already know? Cover:
%   - previous-module content (link the day: \href{/agency/aixi}{AIXI});
%   - technical background, as abilities (doing proofs) and knowledge;
%   - background assumptions or world-views that inform the content;
%   - terminology assumed known.
% Ordinary LaTeX — no special markup.

\begin{itemize}
  \item Comfort with probability on finite sets ($\E$, distributions on a
        finite $\mathcal{X}$).
  \item One variable of calculus: we differentiate $-t \log t$ once.
  \item From \href{/foundations/prerequisites}{Prerequisites (day 0)}: big-$O$
        notation and the idea of a bit.
\end{itemize}

%===============================================================================
% ---- Module intent = the "What you'll learn" box ----------------------------
% Three things: WHAT participants must leave with (if they don't get it, the
% day failed); WHY it matters, connected to the goals of the whole intensive;
% HOW the day gets them there — the story of the day, at a high level.
% Two tests. Completeness: could a participant point a friend at this and say
% "read this and you'll get the idea"? Teachability: does it let a teacher
% short on time decide on the fly what to skip?
% Anti-example: "Students understand entropy." (fails both tests.)
%
% The body is ordinary LaTeX: a single itemize, or groups under
% \subsection*{...} headings (bold subheadings in the box, PDF and web). A
% short sheet may drop the groups and keep one itemize.
\begin{learningoutcomes}
  \subsection*{What}
  \begin{itemize}
    \item Define the entropy of a distribution and compute it by hand for
          small examples.
    \item State and prove Gibbs' inequality, and use it to bound entropy.
    \item Recognise every markup element the \LaTeX{} $\to$ web pipeline
          supports, and copy this file as the starting point for a sheet.
  \end{itemize}

  \subsection*{Why}
  Entropy is the unit of account for everything later in the course that
  measures information: description length, KL divergence, mutual
  information. A participant who cannot manipulate it fluently is stuck from
  day two onwards.

  \subsection*{How}
  A morning lecture states the definitions and proves Gibbs' inequality; an
  exercise session makes the participants reprove it and find the equality
  case themselves; a reading and discussion block connects it to coding.
\end{learningoutcomes}

%===============================================================================
\section{Roadmap for today}
\label{sec:roadmap}
% The time schedule: one line per sub-module, matching the sub-module
% headings under Main content below. Keep the numbers REALISTIC and revise
% them after teaching — a teacher uses this to know if they are on track.
% A "Fast track" for the participant who has an hour, or missed the day, sits
% under Content below.

\begin{itemize}
  \item 10:00 -- 11:00 \hyperref[sec:lecture]{Lecture: entropy}
  \item 11:00 -- 11:15 Break
  \item 11:15 -- 12:30 \hyperref[sec:exercises]{Exercise session}
  \item 12:30 -- 13:30 Lunch
  \item 13:30 -- 15:00 \hyperref[sec:reading]{Reading and discussion}
  \item 15:00 -- 15:30 \hyperref[sec:checkpoint]{Daily checkpoint} and feedback form
\end{itemize}

% Teacher-facing material goes in a teachingnote: a dashed box in the PDF, a
% collapsed drop-down on the web, so it stays out of the participant's reading
% flow. The optional argument is its title (default "Teaching note").
\begin{teachingnote}[Required materials]
Projector and whiteboard. Print the exercise sheet without solutions
(the \texttt{-nosol} PDF the site offers) for anyone who wants paper.
\end{teachingnote}

%===============================================================================
\section{Content}
\label{sec:content}
% Mirrors the Mega Doc: Fast track, then Main content structured into
% sub-modules, then Learn more (which the site puts LAST — see below).

%-------------------------------------------------------------------------------
\subsection{Fast track}
\label{sec:fast-track}
% The route through Main content that hits the most important learning
% outcomes in about an hour — for the participant who missed the day, or the
% one catching up. Point at specific sections, exercises and readings.

To get the core in an hour: read \cref{def:entropy,thm:gibbs} and the proof,
do \cref{ex:warmup}, then skim \cite[ch.~4]{MacKay:03}.

%-------------------------------------------------------------------------------
\subsection{Main content}
\label{sec:main}
% What is taught in the day (6.5 hours, 10:00–18:00), chronological, top to
% bottom, self-contained. Structure it by SUB-MODULE: each distinct activity
% — a lecture, an exercise session, a reading, a discussion, a guest lecture —
% is its own \subsubsection, matching a line of the Roadmap. Every
% sub-module opens with a one- or two-sentence INTENT in a teachingnote: what
% it leads to, how, and its place in the story of the day (same two tests as
% the module intent — it is what lets a teacher decide what to cut).
% Add teachingnotes wherever you know something about teaching THIS thing
% that the next teacher would not think of themselves ("TN, slide 46: ...").
% Future-proofing: exercises need solutions; readings need a persistent
% format; a deck needs both a PDF and an editable source (slides.tex here).
%
% A long exercise sheet can instead stand as its own top-level \section{}s
% after this one; keep the Roadmap pointing at them.

The day runs as three sub-modules — a lecture, an exercise session, and a
reading block — closing with the daily checkpoint.

%...............................................................................
\subsubsection{Lecture: entropy}
\label{sec:lecture}

\begin{teachingnote}[Intent]
Give the participants the definition and one inequality they can prove
themselves, so the afternoon exercises are a rehearsal and not a first
sighting. Everything after the proof of \cref{thm:gibbs} can be cut if time is
short.
\end{teachingnote}

\begin{teachingnote}
TN, slide 12: draw $-t\log t$ on the board before showing the plot; most
participants have never seen that it is concave.
\end{teachingnote}

% ---- Prose, math, cross-references -------------------------------------------
Ordinary prose converts verbatim, including \emph{emphasis}, \textbf{bold},
\texttt{monospace}, and special characters like 50\% or \S3. Inline math is
copied byte-for-byte: $e^{i\pi} + 1 = 0$, and KaTeX handles comparisons like
$x_{<t}$ without trouble. Custom notation defined in the preamble just works:
$\KL{p}{q} \geq 0$, $\argmax_{a} \E[R \mid a]$, for $p, q$ distributions on
$\R^n$. A literal dollar in prose is \$1,000.

Display math works as \[ \sum_{k=1}^{n} k = \frac{n(n+1)}{2}, \]
and as numbered equations you can reference:
\begin{equation}
  \label{eq:basel}
  \sum_{n \geq 1} \frac{1}{n^2} = \frac{\pi^2}{6}.
\end{equation}
Multi-line derivations use \texttt{align}:
\begin{align}
  (a+b)^2 &= (a+b)(a+b) \\
          &= a^2 + 2ab + b^2.
\end{align}
Cross-references resolve to the exact text LaTeX prints: \cref{eq:basel}
(or, number only, \eqref{eq:basel}) sits in \cref{sec:lecture}; below we
meet \cref{def:entropy,thm:gibbs}. Always \verb|\cref|, never a hand-written
``Section 3'' next to a \verb|\ref| — the build flags it.

% ---- Definitions, theorems, proofs -------------------------------------------
% definition / theorem / lemma / proposition / corollary / fact / example share
% one per-section counter. Label anything you want to reference; right after
% \begin{...} is clearest. The optional argument is a name.
\begin{definition}[entropy]
\label{def:entropy}
The \emph{entropy} of a distribution $p$ on a finite set $\mathcal{X}$ is
\[ H(p) = -\sum_{x \in \mathcal{X}} p(x) \log p(x). \]
\end{definition}

\begin{theorem}[Gibbs' inequality]
\label{thm:gibbs}
For any distributions $p, q$ on $\mathcal{X}$,
$-\sum_x p(x) \log q(x) \geq H(p)$, with equality iff $p = q$.
\end{theorem}

\begin{proof}
By definition, $\KL{p}{q} = \sum_x p(x) \log \frac{p(x)}{q(x)} \geq 0$ by
Jensen's inequality applied to the concave function $\log$; rearranging gives
the claim. Equality in Jensen requires $p(x)/q(x)$ constant, i.e.\ $p = q$.
(On the website this proof is collapsible, like a solution.)
\end{proof}

\begin{lemma}
\label{lem:uniform}
$H(p) \leq \log |\mathcal{X}|$, with equality iff $p$ is uniform.
\end{lemma}

\begin{proof}
Apply \cref{thm:gibbs} with $q$ uniform.
\end{proof}

\begin{example}
\label{exa:coin}
A fair coin has $H = \log 2$: one bit. A coin with $p(\text{heads}) = 1$ has
$H = 0$.
\end{example}

% ---- Callouts and remarks ----------------------------------------------------
\begin{callout}[warning]
\label{co:pitfall}
Callouts flag pitfalls, tips, or side notes. Types: \texttt{note} (default),
\texttt{tip}, \texttt{warning}. They may be labelled and referenced, like
this one. Mind the base of the logarithm: bits for $\log_2$, nats for $\ln$.
\end{callout}

\begin{callout}[tip][A titled callout]
A second optional argument titles the box: it replaces the ``Tip'' label on
the PDF frame and heads the box on the web. Handy for a highlighted key
equation, \(H(p) = -\sum_x p(x) \log p(x)\), or a named aside. Keep maths out
of the title itself.
\end{callout}

\begin{remark}[Continuous distributions]
\label{rem:aside}
A \texttt{remark} is an aside; on the website it renders as a note-style
callout. The optional title is appended in parentheses. Referencing works:
see \cref{co:pitfall} above (and this remark is \cref{rem:aside}).
\end{remark}

% ---- Figures ---------------------------------------------------------------
% Two ways. (1) Export to PDF into fig/ and use figure + \includegraphics —
% the build converts it to SVG for the web (.svg/.png in fig/ are served
% as-is). (2) An inline tikzpicture, compiled and rendered for the web
% automatically.
\begin{figure}[ht]
  \centering
  \includegraphics[width=0.45\linewidth]{fig/value-curve.pdf}
  \caption{An example figure from \texttt{fig/}. On the web this becomes a
  \texttt{Figure} component.}
  \label{fig:value}
\end{figure}

See \cref{fig:value}. The same in TikZ, inline:

\begin{center}
\begin{tikzpicture}[node distance=14mm, every node/.style={draw, rounded corners, inner sep=4pt}]
  \node (p) {$p$};
  \node (h) [right=of p] {$H(p)$};
  \draw[->, thick] (p) -- node[draw=none, above, font=\scriptsize] {entropy} (h);
\end{tikzpicture}
\end{center}

% ---- Citations, links, footnotes ---------------------------------------------
Citations come from \texttt{biblo.bib}: the idea of entropy is due to
\cite{Shannon:48}; for a textbook treatment see \cite[ch.~4]{MacKay:03}. On
the web every citation becomes author--year text linking to a References list
at the foot of the page, whatever \verb|\bibliographystyle| the PDF uses.
External links are plain
\href{https://en.wikipedia.org/wiki/Entropy_(information_theory)}{hyperlinks},
and footnotes are supported.\footnote{On the web this becomes a numbered
  footnote: the marker links down to a note at the foot of the page, and the
  note links back. In the PDF it is an ordinary \LaTeX{} footnote.}

%...............................................................................
\subsubsection{Exercise session}
\label{sec:exercises}

\begin{teachingnote}[Intent]
The participants reprove the lecture's inequality with their own hands and
find its equality case, so that ``entropy is maximised by the uniform
distribution'' becomes something they derived rather than heard. Pairs;
\cref{ex:gibbs-hard} is the one to discuss on the board.
\end{teachingnote}

% ---- Exercises, hints, solutions ---------------------------------------------
% \begin{exercise}[Title] then \label{ex:...}; an optional \important right
% after the label stars one of the sheet's key exercises. Subparts are a plain
% enumerate — label an \item to \cref it ("Exercise 3.1(a)"). Every solution
% NAMES its exercise, \begin{solution}[ex:label] (mandatory), so it may sit
% right after the exercise or be collected at the end: the PDF keeps your
% placement, the web always shows each solution under its exercise, and the
% build strips them all for the -nosol downloads. Exercises without solutions
% are allowed; label any exercise a solution or \cref points at.

\begin{exercise}[Entropy warm-up]
\label{ex:warmup}
Let $p$ be a distribution on a finite set $\mathcal{X}$.
\begin{enumerate}
  \item Show that $H(p) \geq 0$.
        \label{ex:warmup-a}
  \item For which $p$ is $H(p) = 0$?
\end{enumerate}
\end{exercise}


\begin{exercise}[Gibbs, the hard direction]
\label{ex:gibbs-hard}
\important
Prove the equality case of \cref{thm:gibbs} directly, without citing Jensen's
equality condition. \note{this one is solved at the end of the sheet instead
of inline, to demonstrate both solution placements.}
\end{exercise}

% Hints have an inline form, \hint{...} -> "[Hint: ...]", and this block form:
% unnumbered, never labelled, rendered where written — a bold lead-in on
% paper, a collapsible drop-down on the web. Hints survive the -nosol build.
\begin{hint}
Consider $f(t) = t - 1 - \log t$: it is nonnegative on $t > 0$ and vanishes
only at $t = 1$.
\end{hint}

\begin{exercise}[Open-ended]
\label{ex:open}
Why might a learning system that compresses its inputs well also predict
them well? No solution is given: this is a discussion exercise.
\end{exercise}

%...............................................................................
\subsubsection{Reading and discussion}
\label{sec:reading}

\begin{teachingnote}[Intent]
Connect the morning's definition to coding, through one well-curated
reading with clear take-aways, then a discussion that compares and
contrasts. Groups of three to four; steer towards the source-coding theorem
if the discussion stalls.
\end{teachingnote}

% A reading session: what to read, how much of it, and what to take away.
% Cite from biblo.bib so the reference lands in the References list.
\begin{itemize}
  \item \cite[ch.~4]{MacKay:03}, sections 4.1--4.3 (about 40 minutes).
        Take-away: the entropy is the average length of the best code.
  \item \cite{Shannon:48}, section 6 only. Take-away: where the axioms for
        $H$ come from.
\end{itemize}

\begin{callout}[note][Discussion prompt]
Two participants disagree: one says entropy is a property of the source, the
other that it is a property of the observer's beliefs about the source. Who
is right, and does \cref{thm:gibbs} settle it?
\end{callout}

%...............................................................................
\subsubsection{Daily checkpoint}
\label{sec:checkpoint}
% Mandatory: a short multiple-choice quiz that is easy for someone who
% followed the day and hard otherwise, aligned with the What above. Link it.

A ten-minute quiz, \href{https://example.com/quiz}{here}, then the daily
feedback form.

% ---- Postprocessing (teacher-facing) -----------------------------------------
% Notes for whoever runs or develops the day next: what to change, what to
% add, what could be explained better, how to make it work with less of you.
% Revise after teaching. Retrospectives about individual participants do not
% belong in the sheet — they go to Leon.
\begin{teachingnote}[Notes for future iterations]
The reading block ran long last time; consider dropping Shannon's section 6
and giving the discussion prompt the extra half hour.
\end{teachingnote}

%-------------------------------------------------------------------------------
\subsection{Learn more}
\label{sec:learn-more}
% Beyond what fits in a day: how to deepen the material, closely related
% work, the research frontier and where to read about it, what someone could
% learn in a week; conferences, seminars, programmes. The build expects this
% to be the LAST section — after all taught content, just before the
% references (or before \appendix if there is one).

\textbf{Deeper.} \cite[ch.~5--6]{MacKay:03} for source coding and the noisy
channel. \textbf{Frontier.} The
\href{https://en.wikipedia.org/wiki/Minimum_description_length}{minimum description
length} literature, which is where this thread is picked up on the
Solomonoff induction day.

%===============================================================================
% ---- Collected solutions (optional) -------------------------------------------
% If you collect solutions at the back of the sheet, wrap the header in
% pdfonly + solutionsonly: the web relocates every solution under its
% exercise and would otherwise show an emptied "Solutions" heading, and the
% -nosol PDF must not show it either. Inside pdfonly you may also \cref this
% appendix from prose the web never sees.
\begin{pdfonly}
\end{pdfonly}


% A solutionsonly block appears ONLY in the with-solutions build: plain
% content here and in the with-solutions PDF, removed entirely from every
% -nosol download (like a solution, but with no box, heading, or exercise
% binding). An answer key, an instructor aside, any spoiler.

% Any \bibliographystyle you like governs the PDF only; the web normalises
% every citation to author--year.
\bibliographystyle{plain}
\bibliography{biblo}

\end{document}
